\documentclass{ctexart}
\usepackage{template}
\usepackage{esint}
\begin{document}\pagestyle{empty}
\begin{center}\Large
北京大学2025-26学年第1学期高数B期末试题
\end{center}
\begin{problem}[1.]
试求如下极限.
\begin{enumerate}[label=\tbf{(\alph*)},leftmargin=*,topsep=0pt,parsep=0pt,itemsep=0pt,partopsep=0pt]
\item \[\lim_{x\to0}\dfrac{1-x^2-\e^{-x^2}}{x^2\sin^2x}\]
\item \[\lim_{(x,y)\to(0,0)}\dfrac{xy(x+y)}{x^2+y^2}\]
\end{enumerate}
\end{problem}
\begin{solution}
\begin{enumerate}[label=\tbf{(\alph*)},leftmargin=*]
\item 利用指数函数在原点的Taylor展开式,有
\[\e^{-x^2}=1-x^2+\dfrac{x^4}{2}+o(x^4)\]
又因为$\sin x/x\to1$,故
\[\lim_{x\to0}\dfrac{1-x^2-\e^{-x^2}}{x^2\sin^2x}
=\lim_{x\to0}\dfrac{-\frac12x^4+o(x^4)}{x^4\left(\frac{\sin x}{x}\right)^2}
=-\dfrac12\]
\item 当$(x,y)\ne(0,0)$时,由$2|xy|\leqslant x^2+y^2$,可得
\[0\leqslant\left|\dfrac{xy(x+y)}{x^2+y^2}\right|
\leqslant\dfrac{|x+y|}{2}
\leqslant\dfrac{|x|+|y|}{2}\]
右端在$(x,y)\to(0,0)$时趋于$0$.由夹逼定理,所求极限为$0$.
\end{enumerate}
\end{solution}
\begin{problem}[2.]
设
\[f(x)=\dfrac{\di^3}{\di x^3}(x^2-1)^3\]
\begin{enumerate}[label=\tbf{(\alph*)},leftmargin=*,topsep=0pt,parsep=0pt,itemsep=0pt,partopsep=0pt]
\item 求$f(x)$的显式表达式.
\item 求$f(x)$的单调区间.
\item 求$f(x)$的极值点.
\item 求$f(x)$的凹凸区间.
\end{enumerate}
\end{problem}
\begin{solution}
\begin{enumerate}[label=\tbf{(\alph*)},leftmargin=*]
\item 先展开被求导的多项式:
\[(x^2-1)^3=x^6-3x^4+3x^2-1\]
连续求三次导数,得到
\[f(x)=120x^3-72x=24x(5x^2-3)\]
\item 对$f$求导,有
\[f'(x)=360x^2-72=72(5x^2-1)\]
因此当$x<-1/\sqrt5$或$x>1/\sqrt5$时,$f'(x)>0$;
当$-1/\sqrt5<x<1/\sqrt5$时,$f'(x)<0$.
故$f$在$(-\infty,-1/\sqrt5)$和$(1/\sqrt5,+\infty)$上严格递增,
在$(-1/\sqrt5,1/\sqrt5)$上严格递减.
\item 在$x=-1/\sqrt5$处,$f'$由正变负,故该点为极大值点,且
\[f\left(-\dfrac1{\sqrt5}\right)=\dfrac{48}{\sqrt5}\]
在$x=1/\sqrt5$处,$f'$由负变正,故该点为极小值点,且
\[f\left(\dfrac1{\sqrt5}\right)=-\dfrac{48}{\sqrt5}\]
\item 再求一次导数,有$f''(x)=720x$.
因此在$(-\infty,0)$上,$f''<0$,曲线为凸弧;
在$(0,+\infty)$上,$f''>0$,曲线为凹弧.
曲线的拐点为$(0,0)$.
\end{enumerate}
\end{solution}
\begin{problem}[3.]
已知直线$L_1,L_2$分别为
\[L_1:x=\dfrac y2=\dfrac z3,\qquad L_2:x=-y=z-1\]
若$L_1\subset\pi$, $L_2\parallel\pi$,求平面$\pi$的方程.
\end{problem}
\begin{solution}
由直线方程可知,$L_1$经过原点,方向向量为$\vect a=(1,2,3)$;
$L_2$的方向向量为$\vect b=(1,-1,1)$.
因为$L_1\subset\pi$且$L_2\parallel\pi$,平面$\pi$的法向量应同时垂直于$\vect a$和$\vect b$.
取
\[\vect n=\vect a\times\vect b=
\begin{vmatrix}
\vect i&\vect j&\vect k\\
1&2&3\\
1&-1&1
\end{vmatrix}=(5,2,-3)\]
又因为平面经过原点,故所求平面方程为
\[5x+2y-3z=0\]
将$L_2$上的点$(s,-s,s+1)$代入左端,恒得$-3\ne0$,
所以$L_2$确与此平面平行且不在平面内.
\end{solution}
\begin{problem}[4.]
设函数
\[f(x,y)=\begin{cases}1,&x^2+y^2\leqslant1,\\x^2+y^2,&x^2+y^2>1,\end{cases}\qquad \alpha\in(0,\pi/2)\]
证明:在点$(\cos\alpha,\sin\alpha)$处,偏导数$f_x(\cos\alpha,\sin\alpha)$与$f_y(\cos\alpha,\sin\alpha)$均不存在.
\end{problem}
\begin{solution}
记$a=\cos\alpha$, $b=\sin\alpha$.由于$0<\alpha<\pi/2$,有$a>0$, $b>0$且$a^2+b^2=1$,从而$f(a,b)=1$.
先沿$x$方向考察差商.当$0<|h|<a$时,
\[(a+h)^2+b^2-1=h(2a+h)\]
其符号与$h$相同.因此
\[\dfrac{f(a+h,b)-f(a,b)}h=
\begin{cases}
0,&-a<h<0,\\
2a+h,&0<h<a.
\end{cases}\]
左、右极限分别为$0$与$2a$,二者不相等,故$f_x(a,b)$不存在.

同理,当$0<|k|<b$时,$a^2+(b+k)^2-1=k(2b+k)$,于是
\[\dfrac{f(a,b+k)-f(a,b)}k=
\begin{cases}
0,&-b<k<0,\\
2b+k,&0<k<b.
\end{cases}\]
左、右极限分别为$0$与$2b$,亦不相等.故$f_y(a,b)$也不存在.
\end{solution}
\begin{problem}[5.]
考虑方程
\[F(x,y,z)=xy+yz+\e^{2xz}-4=0\]
\begin{enumerate}[label=\tbf{(\alph*)},leftmargin=*,topsep=0pt,parsep=0pt,itemsep=0pt,partopsep=0pt]
\item 求一点$M_0(x_0,y_0,z_0)$,使得在$M_0$的某邻域内该方程可唯一确定隐函数$z=z(x,y)$,并验证隐函数存在定理的条件.
\item 设在某解点邻域内由方程确定$z=z(x,y)$,求$\dfrac{\p z}{\p x}$, $\dfrac{\p z}{\p y}$.
\end{enumerate}
\end{problem}
\begin{solution}
\begin{enumerate}[label=\tbf{(\alph*)},leftmargin=*]
\item 取$M_0=(0,1,3)$.有
\[F(0,1,3)=0+3+1-4=0\]
函数$F$在$\R^3$上具有连续的各阶偏导数,并且
\[F_z(x,y,z)=y+2x\e^{2xz},\qquad F_z(0,1,3)=1\ne0\]
故由隐函数存在定理,在$(0,1)$的某邻域内存在唯一的连续可微函数$z=z(x,y)$,
满足$z(0,1)=3$及$F(x,y,z(x,y))=0$,且其图像位于$M_0$的某邻域内.
\item 对恒等式$F(x,y,z(x,y))=0$分别关于$x,y$求偏导,得到
\[\begin{aligned}
y+2z\e^{2xz}+(y+2x\e^{2xz})z_x&=0,\\
x+z+(y+2x\e^{2xz})z_y&=0.
\end{aligned}\]
在$F_z=y+2x\e^{2xz}\ne0$的邻域内,故有
\[\dfrac{\p z}{\p x}=-\dfrac{y+2z\e^{2xz}}{y+2x\e^{2xz}},\qquad
\dfrac{\p z}{\p y}=-\dfrac{x+z}{y+2x\e^{2xz}}\]
式中的$z$均指$z(x,y)$.
\end{enumerate}
\end{solution}
\begin{problem}[6.]
设$z=f(u(x,y),v(x,y))$,其中$u(x,y)=\e^x\cos y$, $v(x,y)=\e^x\sin y$.求$\dfrac{\p^2z}{\p y^2}$.
\end{problem}
\begin{solution}
在$f$具有连续二阶偏导数的条件下,应用复合函数求导法则.
由$u=\e^x\cos y$, $v=\e^x\sin y$,可得
\[u_y=-v,\qquad v_y=u,\qquad u_{yy}=-u,\qquad v_{yy}=-v\]
所以
\[z_y=f_u u_y+f_v v_y=-vf_u+uf_v\]
再关于$y$求导,有
\[\begin{aligned}
z_{yy}
&=f_{uu}u_y^2+2f_{uv}u_yv_y+f_{vv}v_y^2+f_u u_{yy}+f_v v_{yy}\\
&=v^2f_{uu}-2uvf_{uv}+u^2f_{vv}-uf_u-vf_v.
\end{aligned}\]
代回$u,v$的表达式,即得
\[\begin{aligned}
\dfrac{\p^2z}{\p y^2}
={}&\e^{2x}\bigl(\sin^2y\,f_{uu}-2\sin y\cos y\,f_{uv}+\cos^2y\,f_{vv}\bigr)\\
&-\e^x\bigl(\cos y\,f_u+\sin y\,f_v\bigr).
\end{aligned}\]
上式中$f$的各偏导数均在$(u,v)=(\e^x\cos y,\e^x\sin y)$处取值.
\end{solution}
\begin{problem}[7.]
设$F(x,y)=f(x^2+y^2)$,写出$F(x,y)$在$(x_0,y_0)$附近的二阶Taylor展开式,并用佩亚诺余项表示.
\end{problem}
\begin{solution}
记
\[s_0=x_0^2+y_0^2,\qquad h=x-x_0,\qquad k=y-y_0,\qquad \rho=\sqrt{h^2+k^2}\]
为使用二阶Taylor公式,假设$f$在$s_0$附近二阶连续可微.
由复合函数求导法则,在$(x_0,y_0)$处有
\[\begin{aligned}
F_x&=2x_0f'(s_0),&F_y&=2y_0f'(s_0),\\
F_{xx}&=2f'(s_0)+4x_0^2f''(s_0),&
F_{xy}&=4x_0y_0f''(s_0),\\
F_{yy}&=2f'(s_0)+4y_0^2f''(s_0).
\end{aligned}\]
将它们代入二元函数的二阶Taylor公式,整理得到
\[\begin{aligned}
F(x_0+h,y_0+k)
={}&f(s_0)+2f'(s_0)(x_0h+y_0k)\\
&+f'(s_0)(h^2+k^2)+2f''(s_0)(x_0h+y_0k)^2\\
&+o(h^2+k^2),\qquad \rho\to0.
\end{aligned}\]
其中$o(h^2+k^2)$即二阶Peano余项.
\end{solution}
\begin{problem}[8.]
设$f(x,y,z)=x^2-xy+y^2+z^2$,约束条件为平面$x+y+z=1$.求$f$在该平面上的最小值及其取值点,并证明所得结论.
\end{problem}
\begin{solution}
先用Lagrange乘数法寻找可能的极小值点.设
\[\Phi(x,y,z,\lambda)=x^2-xy+y^2+z^2-\lambda(x+y+z-1)\]
约束函数的梯度为$(1,1,1)\ne(0,0,0)$,故条件极值点必须满足
\[\begin{cases}
2x-y=\lambda,\\
-x+2y=\lambda,\\
2z=\lambda,\\
x+y+z=1.
\end{cases}\]
前两式相减,得$x=y$;再代入各式,得到唯一候选点
\[(x,y,z)=\left(\dfrac25,\dfrac25,\dfrac15\right),\qquad f(x,y,z)=\dfrac15\]

为证明这确为整个约束平面上的最小值,令
\[a=x-\dfrac25,\qquad b=y-\dfrac25,\qquad c=z-\dfrac15\]
则$a+b+c=0$.直接展开,得到
\[\begin{aligned}
f(x,y,z)
&=\dfrac15+\dfrac25(a+b+c)+a^2-ab+b^2+c^2\\
&=\dfrac15+\dfrac12(a-b)^2+\dfrac12(a^2+b^2)+c^2\\
&\geqslant\dfrac15.
\end{aligned}\]
等号当且仅当$a=b=c=0$成立.因此$f$在该平面上的最小值为$1/5$,
且仅在$\left(2/5,2/5,1/5\right)$处取得.
\end{solution}
\end{document}
